% !TEX root = ../main.tex
\clearpage
\phantomsection
\section*{5 \textbf{Research Methodology}}
\label{sec:research-methodology}
\addcontentsline{toc}{section}{5 Research Methodology}

This section outlines the overall approach for implementing EndorseRank, comparing it to Adaptive Weighted PageRank (AWP), and evaluating their performance in terms of reputation structure, predictive validity, and computational efficiency.


\phantomsection
\subsection*{5.1 Research Design and Method}
\label{sec:research-design-and-method}
\addcontentsline{toc}{subsection}{5.1 Research Design and Method}

\textbf{Dataset Overview}


This study uses public, on-chain event logs as the primary dataset. The raw data are time-stamped, address-level records emitted by Ethereum smart contracts, which makes them reproducible but also inherently pseudonymous (wallets are identifiers, not verified real-world entities).


\textbf{Core unit of analysis}


Wallet addresses and directed relationships between wallets (owner \(\to\) spender), later represented as a sparse graph.


\textbf{Primary signals}


ERC-20 Approve (and EIP-2612 Permit) authorizations that specify an owner, a spender, and an allowance amount (with token contracts providing the context for each authorization).


\textbf{Sparsity and heavy tails}


Endorsement/authorization graphs are expected to be highly sparse with skewed degree/weight distributions (a small number of contracts/routers receive many approvals), which motivates robust evaluation beyond mean-based summaries.


\textbf{Temporal structure}


Events form a time series; we operationalize EndorseRank on a fixed observation window and construct ``latest state'' edge weights by retaining the most recent relevant allowance information per relationship.


\textbf{Data quality consideration}s


Event logs are generally reliable but can include revocations (amount = 0), ``spend-all by default'' approvals, and occasional indexing gaps; these characteristics are treated explicitly in preprocessing and robustness checks.


Data Collection \& Preprocessing


\textbf{Blockchain Data Source}


Ethereum mainnet (January 1, 2025 -- June 30, 2025) via Google BigQuery or an archive node.


Event Extraction


Approval/permit-derived authorizations: all ERC-20 Approval(owner, spender, value) events, including approvals generated through permit-compatible flows where the token emits the standard Approval event; where permit-specific traces are available, they will be used only to validate the authorization source.


Default Labels


Aave V3 liquidation events will be used as the observable on-chain proxy for default. A wallet will be labelled as defaulted only when the liquidation outcome occurs after the feature-observation window, so that EndorseRank and baseline scores are computed before the outcome they are used to predict.


Sampling


Sampling: the predictive analysis will use the full set of Aave V3-active wallets that appear in the observation window and have sufficient event history for feature construction. Computational benchmarks will use a reproducible random sample of 100,000 wallets with a fixed seed, and robustness checks will compare 50,000- and 200,000-wallet samples to test scalability sensitivity.


\begin{longtable}{>{\raggedright\arraybackslash}p{0.18\linewidth}>{\raggedright\arraybackslash}p{0.39\linewidth}>{\raggedright\arraybackslash}p{0.35\linewidth}}
\toprule
\textbf{Construct} & \textbf{Definition} & \textbf{Associated Variables} \\
\midrule
Reputation\newline Structure & The pattern and relative ordering\newline of wallet reputations produced by\newline a ranking algorithm. & (1) EndorseRank score;\newline AWP score;\newline Rank correlation(Spear- man's \emph{\(\rho\) }/ Kendall's \emph{\(\tau\) }) between EndorseRank and AWP \\
\addlinespace
Risk Predictive\newline Validity & The extent to which a reputation score forecasts actual default events on a DeFi platform. & Default status (binary);\newline EndorseRank's predictive metrics (AUC, Precision/Recall, F1- score);\newline AWP's predictive metrics \\
\addlinespace
Computational\newline Efficiency & The computational resources(time, memory, convergence iterations) required to compute\newline scores. & (7) EndorseRank runtime;\newline AWP runtime;\newline Peak memory usage;\newline Iterations to convergence \\
\addlinespace
\bottomrule
\end{longtable}

\proposalSubheading{1. Graph Construction \& Algorithm Implementation}

\begin{itemize}
\item[\textbullet] \textbf{Approve/Permit Graph }($G_{\text{approve}}$)
\begin{itemize}
\item[\textendash] \textbf{Nodes}: Wallet addresses.
\item[\textendash] Edges: directed owner-to-spender relationships weighted by the latest nonzero allowance from u to v. Because ERC-20 tokens differ in decimals and economic scale, weights will first be token-decimal adjusted; sensitivity checks will compare raw allowance weights, log-transformed weights, and value-normalized weights where reliable price data are available.
\end{itemize}
\item[\textbullet] \textbf{Transfer-Based Graph }($G_{\text{transfer}}$)
\begin{itemize}
\item[\textendash] Same node set; edges weighted by sum of transfer amounts \emph{u }\(\to\) \emph{v}.
\end{itemize}
\item[\textbullet] Algorithms
\end{itemize}
\textbf{EndorseRank: Weighted PageRank on }$G_{\text{approve}}$\textbf{with damping factor d = 0.85. This value is retained because 0.85 is the standard PageRank setting, provides a stable balance between diffusion and teleportation, and enables direct comparison with prior graph-ranking studies.}


\begin{itemize}
\item[\textendash] \textbf{AWP Reproduction}: Adaptive Weighted PageRank on $G_{\text{transfer}}$
\end{itemize}
incorporating time and value decay per Do et al. (2023).


\begin{itemize}
\item[\textbullet] Tools \& Environment
\begin{itemize}
\item[\textendash] Python with NetworkX or GraphFrames; Scipy sparse solvers; Docker containers for reproducibility.
\end{itemize}
\end{itemize}
\proposalSubheading{2. Performance Evaluation Procedures}

\begin{itemize}
\item[\textbullet] Predictive Performance
\begin{itemize}
\item[\textendash] Features: EndorseRank and AWP scores.
\item[\textendash] Labels: Binary default status.
\item[\textendash] Models: Logistic regression and random forest. Logistic regression is used as an interpretable linear baseline, while random forest captures nonlinear interactions and offers a stronger benchmark for mixed on-chain feature patterns.
\item[\textendash] Metrics: AUC, Precision@k, Recall, F1-score.
\end{itemize}
\item[\textbullet] Validation: 5-fold cross-validation, reporting mean +/- SD. Five folds provide a practical bias-variance trade-off while keeping repeated model fitting manageable for multiple baselines, robustness checks, and hardware-limited experiments.
\begin{itemize}
\item[\textendash] Temporal validation will be used alongside cross-validation to prevent look-ahead bias: graph features and reputation scores will be computed from events before the prediction cut-off, while liquidation/default labels will be drawn from a subsequent outcome window. This separation ensures that the model is evaluated on forward-looking predictive performance rather than on information that would not have been available at scoring time.
\end{itemize}
\item[\textbullet] Rank Divergence
\begin{itemize}
\item[\textendash] Compute Spearman's \emph{\(\rho\) }and Kendall's \emph{\(\tau\) }between EndorseRank and AWP.
\item[\textendash] Visualize via scatter plots and top-k Jaccard overlap.
\end{itemize}
\item[\textbullet] Computational Efficiency
\begin{itemize}
\item[\textendash] Hardware: Fixed 22-vCPU, 64 GB RAM server.
\item[\textendash] Measures:
\end{itemize}
\end{itemize}
\(\ast\) Runtime (graph build + PageRank convergence, averaged over 5 runs)


\(\ast\) Peak memory usage


\(\ast\) Iterations to reach tolerance 10\textsuperscript{-}\textsuperscript{6}


\phantomsection
\subsection*{5.1 Data Analysis}
\label{sec:data-analysis}
\addcontentsline{toc}{subsection}{5.1 Data Analysis}

\proposalSubheading{1. Hypothesis Testing}

\begin{itemize}
\item[\textbullet] \textbf{H1 (Divergence)}
\begin{itemize}
\item[\textendash] Test whether rank correlations differ significantly from 1 using a one- sample t-test or Wilcoxon signed-rank test.
\end{itemize}
\item[\textbullet] H2 (Predictive Validity)
\begin{itemize}
\item[\textendash] Compare AUCs via DeLong's test for correlated ROC curves.
\item[\textendash] Compare Precision, Recall, F1-scores via paired bootstrap confidence intervals.
\end{itemize}
\item[\textbullet] H3 (Computational Efficiency)
\begin{itemize}
\item[\textendash] Paired t-tests (or Wilcoxon tests if non-normal) on runtime, memory, and iteration counts.
\end{itemize}
\end{itemize}
\proposalSubheading{2. Robustness Checks}

\begin{itemize}
\item[\textbullet] Vary damping factor \emph{d }\(\in\) \{0\emph{.}75\emph{, }0\emph{.}85\emph{, }0\emph{.}95\}.
\item[\textbullet] Repeat analyses on subgraphs of top ERC-20 tokens by volume.
\item[\textbullet] Assess sensitivity to sample size (50k vs. 200k wallets).
\end{itemize}
\proposalSubheading{3. Reproducibility}

\begin{itemize}
\item[\textbullet] All code, queries, and analysis scripts versioned on GitHub.
\item[\textbullet] Data snapshots and Docker images preserved for full replication of results.
\end{itemize}
This methodological framework ensures a rigorous, reproducible evaluation of EndorseRank's advantages over existing PageRank-based approaches in reputation structure, risk prediction, and computational performance.
